\documentclass[10pt,a4paper]{amsart}
\usepackage{amsmath,amssymb,amsthm,graphicx}
\usepackage{verbatim}
\newenvironment{reference}{\comment}{\endcomment}
%\newenvironment{reference}{}{}
\newenvironment{slogan}{\comment}{\endcomment}
\newenvironment{history}{\comment}{\endcomment}

% For commutative diagrams we use Xy-pic
\usepackage[all]{xy}

% We use 2cell for 2-commutative diagrams.
\xyoption{2cell}
% Theorem environments.
%
\theoremstyle{plain}
\newtheorem{theorem}[subsection]{Theorem}
\newtheorem{proposition}[subsection]{Proposition}
\newtheorem{lemma}[subsection]{Lemma}

\theoremstyle{definition}
\newtheorem{definition}[subsection]{Definition}
\newtheorem{example}[subsection]{Example}
\newtheorem{exercise}[subsection]{Exercise}
\newtheorem{situation}[subsection]{Situation}

\theoremstyle{remark}
\newtheorem{remark}[subsection]{Remark}
\newtheorem{remarks}[subsection]{Remarks}

\numberwithin{equation}{subsection}

% Macros
%
\def\lim{\mathop{\mathrm{lim}}\nolimits}
\def\colim{\mathop{\mathrm{colim}}\nolimits}
\def\Spec{\mathop{\mathrm{Spec}}}
\def\Hom{\mathop{\mathrm{Hom}}\nolimits}
\def\Ext{\mathop{\mathrm{Ext}}\nolimits}
\def\SheafHom{\mathop{\mathcal{H}\!\mathit{om}}\nolimits}
\def\SheafExt{\mathop{\mathcal{E}\!\mathit{xt}}\nolimits}
\def\Sch{\mathit{Sch}}
\def\Mor{\mathop{\mathrm{Mor}}\nolimits}
\def\Ob{\mathop{\mathrm{Ob}}\nolimits}
\def\Sh{\mathop{\mathit{Sh}}\nolimits}
\def\NL{\mathop{N\!L}\nolimits}
\def\CH{\mathop{\mathrm{CH}}\nolimits}
\def\proetale{{pro\text{-}\acute{e}tale}}
\def\etale{{\acute{e}tale}}
\def\QCoh{\mathit{QCoh}}
\def\Ker{\mathop{\mathrm{Ker}}}
\def\Im{\mathop{\mathrm{Im}}}
\def\Coker{\mathop{\mathrm{Coker}}}
\def\Coim{\mathop{\mathrm{Coim}}}

% Boxtimes
%
\DeclareMathSymbol{\boxtimes}{\mathbin}{AMSa}{"02}

%
% Macros for moduli stacks/spaces
%
\def\QCohstack{\mathcal{QC}\!\mathit{oh}}
\def\Cohstack{\mathcal{C}\!\mathit{oh}}
\def\Spacesstack{\mathcal{S}\!\mathit{paces}}
\def\Quotfunctor{\mathrm{Quot}}
\def\Hilbfunctor{\mathrm{Hilb}}
\def\Curvesstack{\mathcal{C}\!\mathit{urves}}
\def\Polarizedstack{\mathcal{P}\!\mathit{olarized}}
\def\Complexesstack{\mathcal{C}\!\mathit{omplexes}}
% \Pic is the operator that assigns to X its picard group, usage \Pic(X)
% \Picardstack_{X/B} denotes the Picard stack of X over B
% \Picardfunctor_{X/B} denotes the Picard functor of X over B
\def\Pic{\mathop{\mathrm{Pic}}\nolimits}
\def\Picardstack{\mathcal{P}\!\mathit{ic}}
\def\Picardfunctor{\mathrm{Pic}}
\def\Deformationcategory{\mathcal{D}\!\mathit{ef}}
\usepackage{cleveref}
\crefname{originallemma}{original lemma}{original lemmas}
\Crefname{originallemma}{Original lemma}{Original lemmas}
\makeatletter
\newlabel{algebra-lemma-completion-faithfully-flat}{{00MC}{source}{}{sourcepointer-1}{}}
\newlabel{algebra-lemma-completion-faithfully-flat@cref}{{[originallemma][0][]00MC}{[1][0][]source}{}{}{}}
\makeatother
\crefname{originallemma}{original lemma}{original lemmas}
\Crefname{originallemma}{Original lemma}{Original lemmas}
\makeatletter
\newlabel{algebra-lemma-descent-regular}{{07NG}{source}{}{sourcepointer-2}{}}
\newlabel{algebra-lemma-descent-regular@cref}{{[originallemma][0][]07NG}{[1][0][]source}{}{}{}}
\makeatother
\crefname{originalreference}{original reference}{original references}
\Crefname{originalreference}{Original reference}{Original references}
\makeatletter
\newlabel{lemma-iso-completions}{{0BG5}{source}{}{sourcepointer-3}{}}
\newlabel{lemma-iso-completions@cref}{{[originalreference][0][]0BG5}{[1][0][]source}{}{}{}}
\makeatother
\crefname{originallemma}{original lemma}{original lemmas}
\Crefname{originallemma}{Original lemma}{Original lemmas}
\makeatletter
\newlabel{more-algebra-lemma-completion-regular}{{07NY}{source}{}{sourcepointer-4}{}}
\newlabel{more-algebra-lemma-completion-regular@cref}{{[originallemma][0][]07NY}{[1][0][]source}{}{}{}}
\makeatother
\crefname{originallemma}{original lemma}{original lemmas}
\Crefname{originallemma}{Original lemma}{Original lemmas}
\makeatletter
\newlabel{morphisms-lemma-base-change-proper}{{01W4}{source}{}{sourcepointer-5}{}}
\newlabel{morphisms-lemma-base-change-proper@cref}{{[originallemma][0][]01W4}{[1][0][]source}{}{}{}}
\makeatother
\crefname{originallemma}{original lemma}{original lemmas}
\Crefname{originallemma}{Original lemma}{Original lemmas}
\makeatletter
\newlabel{properties-lemma-characterize-regular}{{02IT}{source}{}{sourcepointer-6}{}}
\newlabel{properties-lemma-characterize-regular@cref}{{[originallemma][0][]02IT}{[1][0][]source}{}{}{}}
\makeatother
\title{Regularity of the local rings of a scheme over a local ring and of its base change to the completion}
\author{The Stacks Project Authors}
\date{stacks-project:resolve.tex:local-snapshot-20261012; GNU FDL 1.2 or later}
\begin{document}
\maketitle
\noindent\textit{Editorial note.} All delivered mathematical text is the author's own, copied from the pinned Stacks Project chapter file resolve.tex; no statement or proof step was rewritten, reordered or removed.\par\medskip
\noindent\textit{Editorial note.} The counted proof is the author's complete argument: the flat descent of regularity, the comparison of complete local rings at closed fibre points, and the reduction to closed points for the proper case.\par\medskip
\noindent\textit{Editorial note.} The label lemma-port-regularity-to-completion names the counted statement, and the proof environment immediately following it in the source is delivered verbatim as its proof.\par\medskip
\noindent\textit{Editorial note.} The counted statement is carried by the author's own lemma environment in the chapter source.\par\medskip
\noindent\textit{Editorial note.} Pre-open gate: python operations/agent-authoring/tools/precard\_gate.py --source resolve.tex --label lemma-port-regularity-to-completion --start 2996 --end 3026 --local-id 28115 returned RESULT=OK (L28115 free, occupancy FREE, public interval clear). Pre-assembly tag self-check: canonical tag resolved and not present in reference/public-items.jsonl.\par\medskip
\noindent\textit{Editorial note.} The statement and proof are copied from the pinned chapter file; the cited inputs are kept as external pointers with their canonical Stacks tags.\par\medskip
\noindent\textit{Allowed standard prerequisites (editorial).}\begin{itemize}
\item Faithful flatness of the completion. Algebra, the lemma that the completion of a Noetherian local ring is faithfully flat over it. Makes the base change to the completion flat, so that regularity descends.
\item Descent of regularity along a flat morphism. Algebra, the lemma descending regularity of local rings along a flat local homomorphism. Gives the first assertion and one direction of the second.
\item Isomorphism of the complete local rings at the points of the closed fibres. Same chapter, the lemma that the morphism to the base change induces an isomorphism on completions at points over the closed point. Combined with the criterion for regularity through completion, gives the equivalence in the second assertion.
\item Characterization of regularity through closed points and base change of properness. Properties, the characterization of regular schemes by the local rings at closed points, and Morphisms, the stability of properness under base change. Produces the third assertion for a proper morphism.
\end{itemize}\medskip
\ifdefined\hypertarget\hypertarget{sourcepointer-1}{}\fi
\noindent\textit{Source reference.} Original lemma 00MC: sources/stacks-chow/tags.txt, source lines 818--818. Prerequisite pointer invoked at proof line(s) 3014; it names a standard fact of a sibling chapter that the counted argument uses, and every step of the counted proof built on it is retained verbatim in the delivered body.\par\medskip
\ifdefined\hypertarget\hypertarget{sourcepointer-2}{}\fi
\noindent\textit{Source reference.} Original lemma 07NG: sources/stacks-chow/tags.txt, source lines 9761--9761. Prerequisite pointer invoked at proof line(s) 3016; it names a standard fact of a sibling chapter that the counted argument uses, and every step of the counted proof built on it is retained verbatim in the delivered body.\par\medskip
\ifdefined\hypertarget\hypertarget{sourcepointer-3}{}\fi
\noindent\textit{Source reference.} Original reference 0BG5: sources/stacks-resolve/pinned-original-source.tex, source lines 2966--2966. Prerequisite pointer invoked at proof line(s) 3017; it names a standard fact of the same chapter that the counted argument uses, and every step of the counted proof built on it is retained verbatim in the delivered body.\par\medskip
\ifdefined\hypertarget\hypertarget{sourcepointer-4}{}\fi
\noindent\textit{Source reference.} Original lemma 07NY: sources/stacks-chow/tags.txt, source lines 9778--9778. Prerequisite pointer invoked at proof line(s) 3020; it names a standard fact of a sibling chapter that the counted argument uses, and every step of the counted proof built on it is retained verbatim in the delivered body.\par\medskip
\ifdefined\hypertarget\hypertarget{sourcepointer-5}{}\fi
\noindent\textit{Source reference.} Original lemma 01W4: sources/stacks-chow/tags.txt, source lines 2469--2469. Prerequisite pointer invoked at proof line(s) 3022; it names a standard fact of a sibling chapter that the counted argument uses, and every step of the counted proof built on it is retained verbatim in the delivered body.\par\medskip
\ifdefined\hypertarget\hypertarget{sourcepointer-6}{}\fi
\noindent\textit{Source reference.} Original lemma 02IT: sources/stacks-chow/tags.txt, source lines 3312--3312. Prerequisite pointer invoked at proof line(s) 3025; it names a standard fact of a sibling chapter that the counted argument uses, and every step of the counted proof built on it is retained verbatim in the delivered body.\par\medskip
\begin{lemma}
\label{lemma-port-regularity-to-completion}
Let $(A, \mathfrak m, \kappa)$ be a Noetherian local ring.
Let $X \to \Spec(A)$ be a morphism which is locally of finite type.
Set $Y = X \times_{\Spec(A)} \Spec(A^\wedge)$. Let $y \in Y$ with
image $x \in X$. Then
\begin{enumerate}
\item if $\mathcal{O}_{Y, y}$ is regular, then $\mathcal{O}_{X, x}$
is regular,
\item if $y$ is in the closed fibre, then $\mathcal{O}_{Y, y}$ is regular
$\Leftrightarrow \mathcal{O}_{X, x}$ is regular, and
\item If $X$ is proper over $A$, then $X$ is regular
if and only if $Y$ is regular.
\end{enumerate}
\end{lemma}

\begin{proof}
Since $A \to A^\wedge$ is faithfully flat
(Algebra, Lemma \ref{algebra-lemma-completion-faithfully-flat}),
we see that $Y \to X$ is flat. Hence (1) by
Algebra, Lemma \ref{algebra-lemma-descent-regular}.
Lemma \ref{lemma-iso-completions} shows the morphism $Y \to X$
induces an isomorphism on complete local rings at points
of the special fibres. Thus (2) by
More on Algebra, Lemma \ref{more-algebra-lemma-completion-regular}.
If $X$ is proper over $A$, then $Y$ is proper over $A^\wedge$
(Morphisms, Lemma \ref{morphisms-lemma-base-change-proper})
and we see every closed point of $X$ and $Y$ lies in the closed fibre.
Thus we see that $Y$ is a regular scheme if and only if $X$ is so by
Properties, Lemma \ref{properties-lemma-characterize-regular}.
\end{proof}

\end{document}
